\answer{ex:19:1}{(أ)~المشبك الواحد يعطي $6.7\mV$ ($R=66.7\,\text{M}\Omega$)، فبلوغ العتبة أصلًا يحتاج إلى $4$ على الأقل (الثلاثة لا تكفي إلا تقاربيًا)؛ و$8$ في $10\ms$؛ و$14$ في $5\ms$. (ب)~$0.1\ms\ll\tau_m$، وتصحيح التسريب $\sim0.25\%$.}
\answer{ex:19:2}{(أ)~$1.75\cdot10^{10}$ (ربع المازجات يستجيب لكل شيء)، و$4.4\cdot10^9$، و$0.06$ (عند $k=40$ لا يستجيب تقريبًا أبدًا). (ب)~بـ$10^3$ مرة: $2\cdot10^5$ مقابل $200$.}
\answer{ex:19:3}{(أ)~$C(2n,n)=2^{2n-1}$ بتناظر معاملات ذات الحدين. (ب)~$0.93$ و$0.09$؛ عرض الانتقال $\sim\sqrt{2n}$ في $P$، والعرض النسبي $\sim1/\sqrt n$.}
\answer{ex:19:4}{من تغصن واحد $V_{\text{soma}}<\kappa E_E<\theta$ مهما كان عدد المداخل؛ وعند $g_0=0.5g_L$ يلزم $n\ge3$ على كل تغصن.}
\answer{ex:19:5}{$I\ge C\sigma_V/\sigma_t=15\,\text{nA}$، أي $150$ مشبكًا متزامنًا، وهذا غير واقعي؛ أما العصبون الصغير فيكفيه $1\,\text{nA}$، ومع مشبك $4\,\text{nA}$ يكون الارتعاش $5\,\text{\textmu s}$.}
\answer{ex:19:6}{القمة عند الطرف القريب: $0.03\,\tau_m$ و$0.97$ ($L=0.2$)؛ $0.32\,\tau_m$ و$0.66$ ($L=1$)؛ $0.79\,\tau_m$ و$0.33$ ($L=2$). التقدير~\eqref{eq:dendriteload} بالسعة الكاملة لا يصح إلا عند $L\ll1$.}
\answer{ex:19:7}{مسألة مفتوحة. عند $m$ ثابت يزداد زمن الاستجابة خطيًا مع $P$ المحفوظ؛ وعند نسبة ثابتة $m=\varphi n$ يكفّ عن الاعتماد على $n$، فبطء العصبون الكبير نتيجة جمع مداخل كثيرة، لا نتيجة الحجم بحد ذاته.}
